%! Vector-Valued Functions — Unit 4, Lesson 4.9 — Guided Notes
\documentclass[10pt]{article}
\usepackage{precalculus-article}
\usepackage{precalculus-boxes}
\newcommand{\TallMath}[1]{$\displaystyle #1\rule[-1.4em]{0pt}{3.2em}$}

\begin{document}

\pageheader{Unit 4, Lesson 4.9}{Guided Notes}
\namedateperiod

\vspace{0.1in}

\begin{objectivebox}
By the end of this lesson, I will be able to:
\begin{itemize}[leftmargin=1.5em, itemsep=3pt]
  \item \writeline
  \item \writeline
\end{itemize}
\end{objectivebox}

\vspace{0.08in}

\begin{vocabbox}
\termblanklong{vector-valued function}
\termblanklong{position vector}
\termblanklong{velocity vector}
\termblanklong{speed}
\end{vocabbox}

\vspace{0.08in}

\begin{hookbox}
A drone moves so that $x(t)=3t$ and $y(t)=t^{2}-4$ (meters, seconds).
How far from the origin is the drone at $t=2$?  Can a single
mathematical object describe its entire position at once?

\writeline
\end{hookbox}

\vspace{0.08in}

\begin{notesbox}{Section 1 --- Position as a Vector-Valued Function}
\small

\textbf{Key Idea.}\quad Any parametric function $f(t)=(x(t),y(t))$ can be
rewritten as a \textbf{vector-valued function}:
\[
  \vec{p}(t) = \langle \blank{2.5cm},\;\blank{2.5cm}\rangle
             = x(t)\,\hat{\imath} + y(t)\,\hat{\jmath}.
\]

The magnitude of the position vector gives the particle's \textbf{distance from the origin}:
\[
  \|\vec{p}(t)\| = \sqrt{\blank{2.5cm}^{2} + \blank{2.5cm}^{2}}.
\]

\vspace{0.1in}
\textbf{Example.}\quad Drone: $x(t)=3t$, $y(t)=t^{2}-4$.

\begin{tabularx}{\linewidth}{@{}Xl@{}}
Write as a vector-valued function: & $\vec{p}(t) = \langle \blank{3.5cm}\rangle$ \\[12pt]
Find $\vec{p}(2)$: & $\vec{p}(2) = \langle \blank{3.5cm}\rangle$ \\[12pt]
Distance from origin at $t=2$: & $\|\vec{p}(2)\| = \sqrt{\blank{2cm}} = \blank{2cm}$
\end{tabularx}
\end{notesbox}

\vspace{0.08in}

\begin{notesbox}{Section 2 --- Velocity Vector and Speed}
\small

\textbf{Key Idea.}\quad The \textbf{velocity vector} uses the component rates
of change at time $t$:
\[
  \vec{v}(t) = \langle x'(t),\; y'(t)\rangle.
\]

\renewcommand{\arraystretch}{1.5}
\begin{tabular}{lll}
If $x'(t) > 0$: particle moves \blank{2.5cm} &
  \hspace{1.2cm} If $x'(t) < 0$: particle moves \blank{2.5cm} \\
If $y'(t) > 0$: particle moves \blank{2.5cm} &
  \hspace{1.2cm} If $y'(t) < 0$: particle moves \blank{2.5cm}
\end{tabular}

\vspace{0.1in}
\textbf{Speed} = magnitude of velocity = $\|\vec{v}(t)\| = \sqrt{\blank{3cm}}$

\vspace{0.1in}
\textbf{Example (continued).}\quad Drone: $x'(t)=3$, $y'(t)=2t$.

\begin{tabularx}{\linewidth}{@{}Xl@{}}
Velocity at $t=2$: & $\vec{v}(2) = \langle \blank{3.5cm}\rangle$ \\[12pt]
Direction of motion: & $x'(2)=3>0$ $\Rightarrow$ moving \blank{2cm};\quad
  $y'(2)=4>0$ $\Rightarrow$ moving \blank{2cm} \\[12pt]
Speed at $t=2$: & $\|\vec{v}(2)\| = \sqrt{\blank{2.5cm}} = \blank{2cm}$ m/s
\end{tabularx}
\end{notesbox}

\vspace{0.08in}

\begin{notesbox}{Practice}
\small

A particle moves so that $x(t)=t^{2}-4$ and $y(t)=3t$.

\textbf{(a)} Write the position vector: $\vec{p}(t) = \langle \blank{4cm}\rangle$

\vspace{0.1in}
\textbf{(b)} Find $\vec{p}(2)$ and $\|\vec{p}(2)\|$ (exact value).

\vspace{0.06in}
\hspace{1em} $\vec{p}(2) = \langle \blank{3cm}\rangle$ \qquad
$\|\vec{p}(2)\| = \sqrt{\blank{2.5cm}} = \blank{2.5cm}$

\vspace{0.1in}
\textbf{(c)} At $t=2$, with $x'(2)=4$ and $y'(2)=3$: find the speed and describe the direction of motion.

\vspace{0.06in}
\hspace{1em} Speed $= \sqrt{\blank{2.5cm}} = $ \blank{2cm}; \;
direction: \writeline
\end{notesbox}

\end{document}
